\documentclass[11pt]{article}
% ============================================================================
% Single-file manuscript. Preamble and palette are inlined below rather than
% \input from separate files, deliberately.
%
% figs/nm_numbers.tex is the ONE include, and it is generated: make_numbers.py
% writes every \jdg{} and \exd{} macro from the committed run artifacts in
% rna_harness/runs/ceiling/. It is data, not a chapter. No measured number is
% typed into this file by hand, so a printed figure cannot drift from its run.
%
% \jdg{} is measured. \exd{} is expected -- the experiment is declared and has
% not produced the number yet -- and prints with a visible mark so a projection
% can never be read as a measurement.
%
% Build:  python make_numbers.py && tectonic -X compile NM_intervention.tex
% ============================================================================

\usepackage[margin=1in]{geometry}
\usepackage[utf8]{inputenc}
\usepackage[T1]{fontenc}
\usepackage{textcomp}
\usepackage{amsmath,amssymb,amsthm,amsfonts}
\usepackage{booktabs}
\usepackage{graphicx}
\usepackage{xcolor}
\usepackage{array}
\usepackage{enumitem}
\usepackage{placeins}
\usepackage{tabularx}
\usepackage[hidelinks]{hyperref}
\usepackage[numbers,sort&compress]{natbib}
\usepackage{authblk}
\usepackage{caption}
\usepackage{tcolorbox}
\tcbuselibrary{skins,breakable}
\captionsetup{font=small,labelfont=bf,skip=4pt}
\emergencystretch=1.5em
\graphicspath{{figs/}}

\renewcommand{\topfraction}{0.92}
\renewcommand{\bottomfraction}{0.85}
\renewcommand{\textfraction}{0.06}
\renewcommand{\floatpagefraction}{0.72}
\setcounter{topnumber}{4}\setcounter{bottomnumber}{3}\setcounter{totalnumber}{6}

\newtheorem{proposition}{Proposition}
\newtheorem{definition}{Definition}
\newtheorem{assumption}{Assumption}

% ---------------------------------------------------------------- generated numbers
% \jdg{name} prints a measured number defined in figs/nm_numbers.tex. An undefined name is a
% build error rather than a silent blank, because a missing measurement must not print as nothing.
\newcommand{\jdg}[1]{%
  \ifcsname jdg@#1\endcsname\csname jdg@#1\endcsname
  \else\textbf{??#1??}\PackageError{nm}{undefined measured quantity #1}{}\fi}
% \exd{name} prints an EXPECTED number -- declared, not yet produced -- with a visible mark.
\newcommand{\exd}[1]{%
  \ifcsname exd@#1\endcsname\textcolor{S2}{\csname exd@#1\endcsname\,\textsuperscript{\textsc{exp}}}%
  \else\textbf{??#1??}\PackageError{nm}{undefined expected quantity #1}{}\fi}

\newtcolorbox{claimbox}{breakable,enhanced,colback=cbBack,colframe=cbLine,
  boxrule=0.6pt,left=6pt,right=6pt,top=5pt,bottom=5pt,arc=1pt}
\newtcolorbox{scopebox}{breakable,enhanced,colback=sbBack,colframe=sbLine,
  boxrule=0.6pt,left=6pt,right=6pt,top=5pt,bottom=5pt,arc=1pt,
  borderline west={2pt}{0pt}{sbLine}}
\newtcolorbox{withdrawbox}{breakable,enhanced,colback=wbBack,colframe=wbLine,
  boxrule=0.8pt,left=6pt,right=6pt,top=5pt,bottom=5pt,arc=1pt,
  borderline west={3pt}{0pt}{wbLine}}

\definecolor{S1}{HTML}{2A78D6}
\definecolor{S2}{HTML}{EB6834}
\definecolor{S3}{HTML}{1BAF7A}
\definecolor{INK}{HTML}{0B0B0B}
\definecolor{INKB}{HTML}{52514E}
\definecolor{INKC}{HTML}{8A8984}
\definecolor{cbBack}{HTML}{F4F8FD}
\definecolor{cbLine}{HTML}{2A78D6}
\definecolor{sbBack}{HTML}{FBF6F2}
\definecolor{sbLine}{HTML}{EB6834}
\definecolor{wbBack}{HTML}{FDF3F3}
\definecolor{wbLine}{HTML}{B3261E}

% GENERATED by make_numbers.py -- do not edit.
% \jdg{} = measured, from a committed run artifact.
% \exd{} = expected; the experiment is declared but has not produced this number yet.
\expandafter\def\csname jdg@bdNeural\endcsname{23.2\%}
\expandafter\def\csname jdg@bdVienna\endcsname{6.8\%}
\expandafter\def\csname jdg@biasBeta\endcsname{+0.06759}
\expandafter\def\csname jdg@biasCrosscheck\endcsname{\ensuremath{9.71\times10^{-17}}}
\expandafter\def\csname jdg@biasGap\endcsname{\ensuremath{-}0.3487}
\expandafter\def\csname jdg@biasHi\endcsname{\ensuremath{-}0.0206}
\expandafter\def\csname jdg@biasLo\endcsname{\ensuremath{-}0.0239}
\expandafter\def\csname jdg@biasMatchMax\endcsname{0.0021}
\expandafter\def\csname jdg@biasPredicted\endcsname{\ensuremath{-}0.0236}
\expandafter\def\csname jdg@biasRegMax\endcsname{0.0021}
\expandafter\def\csname jdg@calHi\endcsname{+0.5038}
\expandafter\def\csname jdg@calLo\endcsname{+0.2994}
\expandafter\def\csname jdg@calNBins\endcsname{5}
\expandafter\def\csname jdg@calPHi\endcsname{0.9993}
\expandafter\def\csname jdg@calPLo\endcsname{0.0001}
\expandafter\def\csname jdg@clClean\endcsname{3993}
\expandafter\def\csname jdg@clDesigns\endcsname{158}
\expandafter\def\csname jdg@clDirty\endcsname{324}
\expandafter\def\csname jdg@clLeaked\endcsname{14}
\expandafter\def\csname jdg@clLeakedPct\endcsname{8.9\%}
\expandafter\def\csname jdg@clNeuralAll\endcsname{+0.20692}
\expandafter\def\csname jdg@clNeuralClean\endcsname{+0.21549}
\expandafter\def\csname jdg@clThr\endcsname{0.80}
\expandafter\def\csname jdg@clViennaAll\endcsname{+0.12059}
\expandafter\def\csname jdg@clViennaClean\endcsname{+0.12881}
\expandafter\def\csname jdg@cvAlt\endcsname{+0.0891}
\expandafter\def\csname jdg@cvAltHitAU\endcsname{0.406}
\expandafter\def\csname jdg@cvAltHitCG\endcsname{0.628}
\expandafter\def\csname jdg@cvAltWeight\endcsname{1982}
\expandafter\def\csname jdg@cvConv\endcsname{+0.0198}
\expandafter\def\csname jdg@cvConvP\endcsname{0.0199}
\expandafter\def\csname jdg@cvDesigns\endcsname{158}
\expandafter\def\csname jdg@cvHitAU\endcsname{0.908}
\expandafter\def\csname jdg@cvHitCG\endcsname{0.898}
\expandafter\def\csname jdg@cvHitGU\endcsname{0.693}
\expandafter\def\csname jdg@cvIntv\endcsname{+0.0527}
\expandafter\def\csname jdg@cvIntvP\endcsname{0.0025}
\expandafter\def\csname jdg@cvPaired\endcsname{+0.0329}
\expandafter\def\csname jdg@cvPairedCi\endcsname{[+0.0024,\,+0.0638]}
\expandafter\def\csname jdg@cvPairs\endcsname{4293}
\expandafter\def\csname jdg@cvRatio\endcsname{0.38}
\expandafter\def\csname jdg@cvStrata\endcsname{549}
\expandafter\def\csname jdg@cvSwing\endcsname{4.50}
\expandafter\def\csname jdg@cvWeight\endcsname{892}
\expandafter\def\csname jdg@gmNeural\endcsname{+0.20692}
\expandafter\def\csname jdg@gmVienna\endcsname{+0.12059}
\expandafter\def\csname jdg@lkDeposits\endcsname{14}
\expandafter\def\csname jdg@lkExact\endcsname{0}
\expandafter\def\csname jdg@lkGt08\endcsname{6.37\%}
\expandafter\def\csname jdg@lkGt09\endcsname{3.84\%}
\expandafter\def\csname jdg@lkOurs\endcsname{14416}
\expandafter\def\csname jdg@lkTrain\endcsname{861{,}916}
\expandafter\def\csname jdg@lkVerbatim\endcsname{264}
\expandafter\def\csname jdg@mcNeuralCorr\endcsname{\ensuremath{-}0.7019}
\expandafter\def\csname jdg@mcNeuralFull\endcsname{131.4\%}
\expandafter\def\csname jdg@mcNeuralMae\endcsname{0.4743}
\expandafter\def\csname jdg@mcNeuralZero\endcsname{103.2\%}
\expandafter\def\csname jdg@mcSep\endcsname{0.053}
\expandafter\def\csname jdg@mcViennaCorr\endcsname{+0.7026}
\expandafter\def\csname jdg@mcViennaFull\endcsname{-842.5\%}
\expandafter\def\csname jdg@mcViennaMae\endcsname{1.7358}
\expandafter\def\csname jdg@nAU\endcsname{1116}
\expandafter\def\csname jdg@nAdmitted\endcsname{4317}
\expandafter\def\csname jdg@nCG\endcsname{3061}
\expandafter\def\csname jdg@nConstructs\endcsname{13702}
\expandafter\def\csname jdg@nDropped\endcsname{0}
\expandafter\def\csname jdg@nExclNoOutcome\endcsname{27}
\expandafter\def\csname jdg@nExclNonCanon\endcsname{73}
\expandafter\def\csname jdg@nExclTwoPos\endcsname{9125}
\expandafter\def\csname jdg@nGU\endcsname{140}
\expandafter\def\csname jdg@nUniqueSeq\endcsname{4475}
\expandafter\def\csname jdg@nnAsym\endcsname{\ensuremath{0}}
\expandafter\def\csname jdg@nnCorrP\endcsname{+0.4879}
\expandafter\def\csname jdg@nnCorrS\endcsname{+0.5752}
\expandafter\def\csname jdg@nnOffTarget\endcsname{0.0026}
\expandafter\def\csname jdg@nnOnTarget\endcsname{0.7812}
\expandafter\def\csname jdg@nnRatio\endcsname{300}
\expandafter\def\csname jdg@panCorrHi\endcsname{0.92}
\expandafter\def\csname jdg@panCorrLo\endcsname{0.75}
\expandafter\def\csname jdg@panHi\endcsname{+0.1311}
\expandafter\def\csname jdg@panLo\endcsname{+0.1056}
\expandafter\def\csname jdg@panN\endcsname{4}
\expandafter\def\csname jdg@panNames\endcsname{Turner2004, Turner1999, Andronescu2007, Langdon2018}
\expandafter\def\csname jdg@resAgree\endcsname{\ensuremath{2\times10^{-17}}}
\expandafter\def\csname jdg@resCells\endcsname{144}
\expandafter\def\csname jdg@resCellsNeg\endcsname{91}
\expandafter\def\csname jdg@resCellsP\endcsname{0.0019}
\expandafter\def\csname jdg@resNeural\endcsname{+0.0605}
\expandafter\def\csname jdg@resNeuralCi\endcsname{[+0.0329,\,+0.0956]}
\expandafter\def\csname jdg@resNeuralP\endcsname{<0.001}
\expandafter\def\csname jdg@resOut\endcsname{+0.0539}
\expandafter\def\csname jdg@resOutCi\endcsname{[+0.0220,\,+0.0911]}
\expandafter\def\csname jdg@resOutP\endcsname{<0.001}
\expandafter\def\csname jdg@resStrata\endcsname{8}
\expandafter\def\csname jdg@resStrataNeg\endcsname{6}
\expandafter\def\csname jdg@resVienna\endcsname{\ensuremath{-}0.0304}
\expandafter\def\csname jdg@resViennaCi\endcsname{[\ensuremath{-}0.0628,\,+0.0105]}
\expandafter\def\csname jdg@resViennaP\endcsname{0.979}
\expandafter\def\csname jdg@resWeight\endcsname{902}
\expandafter\def\csname jdg@rpCorpusMean\endcsname{0.4124}
\expandafter\def\csname jdg@rpCorpusSd\endcsname{0.2534}
\expandafter\def\csname jdg@rpMaxDiff\endcsname{0.0014}
\expandafter\def\csname jdg@rpMedDiff\endcsname{0.000081}
\expandafter\def\csname jdg@rpN\endcsname{4688}
\expandafter\def\csname jdg@rpPearson\endcsname{1.0000}
\expandafter\def\csname jdg@rpSpearman\endcsname{1.0000}
\expandafter\def\csname jdg@scDeposits\endcsname{1024}
\expandafter\def\csname jdg@scFlips\endcsname{8{,}658}
\expandafter\def\csname jdg@scHits\endcsname{11}
\expandafter\def\csname jdg@skNeural\endcsname{+0.20692}
\expandafter\def\csname jdg@skNeuralP\endcsname{0.002}
\expandafter\def\csname jdg@skNeuralWt\endcsname{+0.24929}
\expandafter\def\csname jdg@skVienna\endcsname{+0.12059}
\expandafter\def\csname jdg@skViennaP\endcsname{0.002}
\expandafter\def\csname exd@xNatFamilies\endcsname{3}
\expandafter\def\csname exd@xNatFlips\endcsname{3954}
\expandafter\def\csname exd@xNatProbes\endcsname{2}
\expandafter\def\csname exd@xNatSkill\endcsname{positive}
\expandafter\def\csname exd@xNatUsable\endcsname{160}


\title{\bf Scoring RNA Structure Prediction on Interventions\\That Have Already Been Performed\\[4pt]
\large Compensatory mutagenesis in a chemical-probing corpus yields \jdg{nAdmitted} controlled
two-nucleotide\\experiments; scoring models on them separates two model classes that\\
agreement with a reference structure cannot}

\author{Wenhui Chen\thanks{\texttt{wenhui.chen@titanresearchlabs.ai}}}
\date{}

\begin{document}
\maketitle

\begin{abstract}
\noindent
An RNA structure benchmark asks whether a predicted structure agrees with a reference. On the
corpus studied here that question has almost no resolution left: the inferred structure already
recovers the pairs under comparison for \jdg{cvHitCG} of G--C targets and \jdg{cvHitAU} of A--U
targets, and its answer moves by a factor of \jdg{cvSwing} depending on which pairs are admitted
to the comparison. We show that the same corpus already contains a different and better question.

Compensatory double mutants are controlled two-nucleotide interventions: each is byte-identical to
its own wild type outside two positions, and those two positions exchange their bases, so the
unordered pair type is preserved while the stacking context is not. We give an eligibility
criterion that recovers them, \jdg{nAdmitted} of them from \jdg{nConstructs} constructs, together
with the counted reasons every other construct is refused, and a tool that reproduces the numbers
in this paper from the public corpus file. Because these interventions were performed and
measured, a structure predictor can be \emph{scored} on them rather than adjusted for: we compare
what a model says a two-nucleotide substitution should do against what the measurement says it
did.

The instrument is calibrated before it is used. Nearest-neighbour thermodynamics has real skill
here --- measured rescue rises monotonically from \jdg{calLo} to \jdg{calHi} across bins of the
model's own equilibrium pairing probability, within-stratum slope \jdg{skVienna} ($p=\jdg{skViennaP}$)
--- and a transformer trained on chemical mapping has more, \jdg{skNeural} ($p=\jdg{skNeuralP}$),
while agreeing with the thermodynamic model only at Pearson \jdg{nnCorrP}. Neither was trained on
the mutants, whose library postdates the neural weights; a four-level leakage audit against
\jdg{lkTrain} sequences from the neural model's own source libraries finds \jdg{lkExact} exact
matches and \jdg{lkGt08} of design windows at $\ge0.8$ identity, and excising the affected
\jdg{clLeaked} of \jdg{clDesigns} design cells leaves both skills unchanged.

What the instrument then reads out is a separation between the two model classes. On the same
\jdg{resWeight} matched pair-weights, the measurement separates G--C from A--U targets by
\jdg{resOut} (\jdg{resOutCi}); the neural model's own prediction separates them by \jdg{resNeural}
(\jdg{resNeuralCi}), the same direction and comparable magnitude; the thermodynamic model's
prediction separates them by \jdg{resVienna} (\jdg{resViennaCi}), which does not resolve from
zero. \textbf{The instrument distinguishes something one model class represents and the other does
not}, and it does so on molecules that were synthesised and probed, with no reference structure
entering the comparison.

We also withdraw a metric. Earlier versions of this work reported the fraction of the identity
contrast surviving after the model's prediction is held fixed. Calibrating that estimator against
a known mediated fraction --- real predictor as mediator, synthetic outcome --- shows it does not
track the truth at this corpus's signal-to-noise (Pearson \jdg{mcNeuralCorr} for the neural arm,
mean absolute error \jdg{mcNeuralMae}). Those numbers are retracted in
Section~\ref{sec:withdrawal} rather than in an appendix, and the claim above deliberately uses
none of them.
\end{abstract}

% =====================================================================
\section{Introduction}
\label{sec:intro}

Every benchmark that scores an RNA secondary-structure predictor asks a version of the same
question: does the structure the model produced match a structure someone else believes to be
right \citep{mathews2004incorporating,hajdin2013accurate,waymentsteele2022rna}. The question is
natural, it has organised the field for thirty years, and on any corpus where models have become
good it eventually runs out of room. Two things go wrong at once. The first is saturation: when
the pairs being compared are already recovered nine times out of ten, the remaining variation is
mostly noise about which structure a degenerate profile happened to select. The second is a
degree of freedom that is rarely reported --- \emph{which} pairs enter the comparison. Restricting
to the designed target pairs, or admitting every pair a predictor emits, are both defensible
choices, and on the corpus below they move the same measurement by a factor of \jdg{cvSwing}.

There is a different question available, and the data to ask it already exists. Chemical-probing
campaigns that include mutant libraries have, for more than a decade, synthesised and measured
molecules that differ from a wild type at a small number of positions
\citep{kladwang2011two,cordero2015rich,tian2014correcting}. Among those mutants is a class with an
unusual property. A \emph{compensatory} double mutant of a target base pair exchanges the two
bases: a G--C becomes a C--G, an A--U becomes a U--A. Outside those two positions the molecule is
unchanged; the unordered pair type is unchanged; what changes is the nearest-neighbour context and
whatever else the exchange disturbs. That is a controlled two-nucleotide experiment with its own
internal control, it has been performed, and its result has been recorded.

We use those interventions to score models rather than to infer structures. This is the point at
which our use departs from the technique's own literature. Mutate-and-map and mutate-map-rescue
use mutations to \emph{establish} which nucleotides pair \citep{kladwang2011two,tian2014correcting,
cheng2017rna}; the mutation is an instrument aimed at the molecule. We aim it at the model instead:
given a mutant that was actually made, what did the model predict would happen to the target pair,
and did the measurement agree? No reference structure enters that comparison, so the degeneracy
that makes profile-based scoring awkward does not arise, and the saturation does not either,
because the quantity being predicted is a change rather than a state.

Three things had to be true for this to be a method rather than an observation, and the paper is
organised around establishing them. The interventions must be recoverable by a stated rule rather
than by inspection, which is Section~\ref{sec:mining}. The instrument must be shown to have
resolution before any conclusion is drawn from it, which is Section~\ref{sec:scoring} --- and the
predictors we score are not innocent of the corpus, so the leakage audit lives there too, in the
main text. And the readout must produce something the incumbent readout cannot, which is
Section~\ref{sec:result}.

It also has to be possible to be wrong in public. Section~\ref{sec:withdrawal} retracts a
quantity this work previously reported --- the fraction of the identity contrast surviving after a
model's prediction is held fixed --- because calibrating that estimator against a known truth
shows it does not measure what it was read as measuring. The retraction is placed in the results,
next to the calibration table that forced it, and the paper's claims are stated so that none of
them depends on it.

\paragraph{Contributions.}
\begin{enumerate}[leftmargin=1.4em,itemsep=2pt]
\item An eligibility criterion that recovers controlled two-nucleotide interventions from a
probing corpus, with named and counted exclusions, ambiguity refused rather than resolved, and no
threshold that can be tuned to raise the yield. It admits \jdg{nAdmitted} interventions here and
refuses \jdg{nExclTwoPos}, \jdg{nExclNonCanon} and \jdg{nExclNoOutcome} constructs for three
stated reasons.
\item A demonstration that the incumbent readout is exhausted on this corpus and carries a
researcher degree of freedom worth \jdg{cvSwing}$\times$, measured against the interventional
readout on the same molecules, the same pairs and the same strata.
\item Three instrument readouts --- sensitivity, calibration and residual --- with the first two
established before the third is reported, and a four-level leakage audit of the one predictor that
could plausibly have memorised these molecules.
\item The result: a separation between a thermodynamic and a neural predictor on the identity axis
that the conventional readout cannot express, recomputed by two independently written estimators
agreeing to \jdg{resAgree}.
\item \texttt{designaudit}, which performs all of the above on a corpus file and whose selftest
reproduces this paper's headline numbers from the public data.
\end{enumerate}

% =====================================================================
\section{Background}
\label{sec:background}

\paragraph{What probing measures, and what it does not.}
Reagents that acylate flexible riboses or methylate unpaired Watson--Crick faces report per
nucleotide on accessibility \citep{merino2005rna,wilkinson2006selective,weeks2010advances,
siegfried2014rna}, and mutational profiling made the readout cheap enough to apply to libraries
rather than to single molecules \citep{zubradt2017dms,marinus2021novel}. A reactivity profile is
not a list of base pairs, though, and the standard route from one to the other is to use the
profile as a constraint on a folding algorithm \citep{deigan2009accurate,hajdin2013accurate}. That
map is one-to-many. Two structures consistent with one profile cannot be separated by the profile,
which is why a benchmark built on profile-derived structures inherits a degeneracy it cannot see.

\paragraph{Covariation and compensatory mutagenesis.}
The argument that survives that objection is older than the models. Compensatory changes across
homologues established the base pairs of ribosomal RNA before any structure was solved
\citep{gutell1994lessons,michel1990modelling}, and compensatory mutagenesis is the same argument
performed rather than observed: disrupt a putative pair with one mutation, disrupt it again with
the partner mutation, and ask whether the pair of mutations together restores what either alone
destroyed. Mutate-and-map applied this systematically \citep{kladwang2011two,cordero2015rich}, and
mutate-map-rescue added the compensatory arm that this paper mines
\citep{tian2014correcting,cheng2017rna}. In every one of those studies the target of inference is
the molecule. We keep the experiment and change the target of inference to the model.

\paragraph{Benchmarks in the neighbourhood.}
Three efforts are close enough to be worth separating from ours explicitly. EternaBench scored
structure packages against a large body of experimental data and used the result to train an
improved model \citep{waymentsteele2022rna}; its question is agreement, ours is response.
RNAGym assembles deep mutational scanning assays and reactivity profiles into a common benchmark
\citep{rnagym2025}, and keeps variant-effect prediction and structure prediction as \emph{separate}
tasks --- the fitness axis scores sequence models on a functional readout, the structure axis
scores predictors on agreement. Our primitive joins them: a structure model is scored by whether
its structural prediction responds correctly to a mutation. And the mutate-and-map literature owns
the experiment itself, including a public server that will predict a mutate-and-map profile
\citep{ribonanzanet2024}; what it does not do is treat the resulting predictions as a benchmark
with a control for who chose the pairs.

\paragraph{Corpora made by design programs.}
The molecules here were produced by design methods competing on shared targets, and a corpus
assembled that way carries its producers' preferences as covariates. This hazard is recognised
qualitatively in several fields --- feedback in recommender systems \citep{chaney2018algorithmic},
generator--judge contamination in language-model evaluation \citep{li2025preference}, characteristic
biases of protein design programs \citep{protdbench2026} --- and the correction it demands here is
quantitative, so we state it as an assumption in Section~\ref{sec:method} and estimate its size in
Section~\ref{sec:conventional}. The estimator theory behind that correction is developed in a
companion study; nothing in this paper depends on reading it.

% =====================================================================
\section{Method}
\label{sec:method}

\subsection{The eligibility criterion}
\label{sec:mining}

Let a corpus supply, for each construct, a sequence, a design window $[\ell,u)$ inside which the
design was written, the target structure the designer was given, a wild-type construct for the
same design cell, and a measured outcome. Write $P$ for the set of pairs in the declared target
structure and $w$, $m$ for the wild-type and mutant sequences.

\begin{definition}[Orientation flip]
\label{def:flip}
A construct is an \emph{admissible intervention} on the pair $(i,j)\in P$ when all five of the
following hold:
\begin{enumerate}[leftmargin=1.6em,itemsep=1pt,label=(\roman*)]
\item $|w| = |m|$;
\item $\{k \in [\ell,u) : w_k \neq m_k\} = \{i,j\}$ exactly --- two positions, both inside the
design window, and no others;
\item $(i,j) \in P$;
\item $\{w_i, w_j\} \in \{\mathrm{A\!-\!U}, \mathrm{C\!-\!G}, \mathrm{G\!-\!U}\}$;
\item $m_i = w_j$ and $m_j = w_i$.
\end{enumerate}
\end{definition}

Clause (v) is what makes the intervention controlled. The unordered pair type is preserved, so the
comparison is not between a pair and its absence; what the substitution changes is the
nearest-neighbour context, the identity of the base presented to each neighbour, and whatever
tertiary contacts depended on the orientation. Clause (ii) is what makes it two nucleotides rather
than an unknown number: constructs in these libraries carry adapters, padding and barcodes outside
the design window, and two constructs differing only by a compensatory flip inside the window can
differ at a dozen positions outside it.

Everything else is refused under a named reason and counted. A mutant admitting more than one
wild type, or whose two changed positions are explained by more than one declared pair, is
excluded as \emph{ambiguous} rather than resolved by a rule --- resolving it would let the corpus
choose which intervention it supplies. Corpus quality fields (signal-to-noise, read depth) are
carried on each admitted intervention and reported, and they never enter the criterion.

\begin{scopebox}
\textbf{What an orientation flip cannot say.} It preserves the unordered pair type. It therefore
does not deliver the causal effect of \emph{being} a G--C rather than an A--U; it delivers whether
a model's prediction of the consequences of this substitution was borne out. Every claim below is
about the second quantity. This limit belongs beside the numbers and not in a closing paragraph.
\end{scopebox}

\subsection{The outcome contains no model}
\label{sec:outcome}

The outcome is the corpus's rescue factor, and because everything in this paper is a statement
about it we verified how it is produced rather than trusting the column. Writing $d_i$, $d_j$ and
$d_{ij}$ for the reactivity profiles of the two single mutants and the double mutant, each minus
the wild-type profile and restricted to the design window,
\begin{equation}
R \;=\; 1 \;-\; \frac{\lVert d_{ij}\rVert}{\sqrt{\lVert d_i\rVert^2 + \lVert d_j\rVert^2}}\,.
\label{eq:rescue}
\end{equation}
Recomputing \eqref{eq:rescue} from the raw per-nucleotide reactivity columns for all \jdg{rpN}
double mutants whose two single mutants and wild type are all present reproduces the corpus's own
value at Pearson $\jdg{rpPearson}$ and Spearman $\jdg{rpSpearman}$ --- same mean
\jdg{rpCorpusMean} and same standard deviation \jdg{rpCorpusSd} --- with a maximum absolute
deviation of \jdg{rpMaxDiff} and a median of \jdg{rpMedDiff}. \textbf{The outcome is measurement-derived end to end; no folding model appears
anywhere in it.} Nothing in this paper is therefore circular in the way a profile-constrained
structure comparison would be.

\subsection{Estimation, and the one assumption it needs}
\label{sec:estimation}

Contrasts are computed inside the design. Write $s(r)$ for a stratum --- the design cell, helix
depth capped at three, and whether the pair is crossed by another --- and compare pair types
within it, weighting each stratum by $\min(n_A(s), n_B(s))$ and pooling. Inference permutes the
\emph{outcome} inside the fixed stratification; intervals resample the puzzle, which is the
independent unit here because design cells that share a target are not independent of one another.

\begin{assumption}[Placement is the confounder that matters]
\label{ass:placement}
Conditional on the design cell, helix depth and crossing, the pair type is as good as randomly
assigned with respect to the outcome.
\end{assumption}

Assumption~\ref{ass:placement} is needed because the molecules were not randomised: design
programs choose both what a pair is and where it goes, and the two choices are correlated. A
controlled-truth benchmark quantifies the price of ignoring it. On synthetic data matched to this
corpus's own depth response $\beta = \jdg{biasBeta}$ and its own pair-type depth gap
$g = \jdg{biasGap}$, an estimator that ignores placement is biased by $\beta g = \jdg{biasPredicted}$,
and the measured bias across cases runs \jdg{biasLo} to \jdg{biasHi}. Regression adjustment and
matching both remove it --- worst-case absolute bias \jdg{biasRegMax} and \jdg{biasMatchMax}
respectively --- so the requirement is to \emph{condition on placement}, not to use any particular
estimator. We say this explicitly because the natural overstatement is that matching is required,
and on this design it is not. A vectorised and a readable implementation of all four estimators
agree to \jdg{biasCrosscheck}.

% =====================================================================
\section{The conventional readout on this corpus}
\label{sec:conventional}

Before proposing a readout it is worth showing what the incumbent one does here, on the same
molecules and with everything else held equal. We take the \jdg{cvDesigns} designs for which both
readouts are defined, intersect to the \jdg{cvPairs} pairs both can score, and verify in code that
the two analyses receive byte-identical stratum weights (\jdg{cvWeight} pair-weights over
\jdg{cvStrata} strata) before either is computed.

\paragraph{It is saturated.} The structure inferred from the wild-type profile already contains
the target pair for \jdg{cvHitCG} of G--C targets, \jdg{cvHitAU} of A--U targets and
\jdg{cvHitGU} of G--U targets. Whatever a benchmark built on that agreement can still resolve, it
is resolving in the last tenth of the range.

\paragraph{It gives a smaller answer on the same data.} The identity contrast read
interventionally is \jdg{cvIntv} ($p=\jdg{cvIntvP}$); read conventionally, from whether the
inferred structure contains the pair, it is \jdg{cvConv} ($p=\jdg{cvConvP}$), a ratio of
\jdg{cvRatio}. Because the two readouts are defined on the same molecules and the same pairs, the
difference can be tested directly: a design-clustered bootstrap of the paired difference gives
\jdg{cvPaired} (\jdg{cvPairedCi}), which excludes zero. The two readouts are measuring different
things, and it is not a matter of one being noisier.

\paragraph{It has a knob.} Scoring instead every target pair of the same designs --- a choice no
less defensible, and one that no longer satisfies the matched-pair-set constraint --- returns
\jdg{cvAlt} on \jdg{cvAltWeight} pair-weights, a factor of \jdg{cvSwing} away from
\jdg{cvConv}. The hit rates on that larger set are \jdg{cvAltHitCG} and \jdg{cvAltHitAU}, so the
enlargement is not benign: it adds pairs the inferred structure mostly does not contain. The
interventional readout has no corresponding freedom, because it is defined only where a
compensatory mutant was made, and which mutants were made is a fact about the corpus rather than a
choice by the analyst.

\begin{claimbox}
\textbf{Claim 1.} On this corpus the conventional structure-agreement readout is near saturation
for the pairs under comparison, and its answer moves by \jdg{cvSwing}$\times$ with a defensible
change in which pairs are admitted. The interventional readout, on the same molecules and strata,
returns a contrast \jdg{cvPaired} (\jdg{cvPairedCi}) larger and admits no such choice.
\end{claimbox}

% =====================================================================
\section{Recovering the interventions}
\label{sec:mining-results}

Applying Definition~\ref{def:flip} to the corpus's \jdg{nConstructs} constructs admits
\jdg{nAdmitted} interventions --- \jdg{nCG} G--C, \jdg{nAU} A--U and \jdg{nGU} G--U --- and
refuses the rest for three reasons: \jdg{nExclTwoPos} constructs do not differ from their wild
type at exactly two in-window positions, \jdg{nExclNonCanon} carry a non-canonical wild-type pair,
and \jdg{nExclNoOutcome} have no measured outcome. Clause (ii) of Definition~\ref{def:flip} is
asserted in code on every admitted record, and \jdg{nDropped} records are dropped by that
assertion, which is the check working rather than a coincidence: the corpus's own annotation and
the sequences agree.

The \jdg{nAdmitted} interventions span \jdg{nUniqueSeq} distinct sequences once wild types are
counted. Every one of them is a molecule that was synthesised and probed.

\paragraph{The tool.} \texttt{designaudit} implements the criterion, the three instrument readouts
of Section~\ref{sec:scoring}, and the estimation of Section~\ref{sec:estimation}, as
\texttt{mine}, \texttt{score} and \texttt{contrast}. Its \texttt{selftest} re-derives this paper's
headline numbers from the public corpus file, independently of the scripts that produced them,
and prints every check whether it passes or fails; a declared check that cannot run is reported as
a failure rather than skipped. Four rules are enforced as code rather than documented: outcome
permutation under a fixed design, same-population denominators, a placebo mediator that always
runs, and a linearity test before any product-of-coefficients quantity is reported.

% =====================================================================
\section{Scoring predictors}
\label{sec:scoring}

A predictor enters as a callable that returns, for a sequence and a position pair, a number it
believes about that pair forming. It need not be a probability, and it need not be thermodynamic;
it must only be comparable across the sequences of one corpus. Three readouts come back, and the
order matters.

\subsection{Sensitivity, then calibration}

\paragraph{The thermodynamic model has skill here.} Taking the equilibrium probability that the
flipped pair forms in the mutant's own ensemble, measured rescue rises monotonically across bins of
that probability, from \jdg{calLo} in the lowest of \jdg{calNBins} bins to \jdg{calHi} in the
highest, with the model's probability running \jdg{calPLo} to \jdg{calPHi} across the same bins.
The within-stratum slope is \jdg{skVienna} ($p=\jdg{skViennaP}$). This is stated first because
without it nothing downstream could be read as being about the model: a model that predicted
nothing could not be shown to be missing anything.

\paragraph{A model with no thermodynamics in it has more.} RibonanzaNet's secondary-structure head
is a linear map on symmetrised pairwise features of a transformer trained on roughly two million
chemical-mapping profiles; it contains no nearest-neighbour parameter table, no partition function
and no ensemble in the thermodynamic sense. Its within-stratum slope is \jdg{skNeural}
($p=\jdg{skNeuralP}$), higher than the thermodynamic model's, and the two predictors agree with
each other only at Pearson \jdg{nnCorrP} and Spearman \jdg{nnCorrS} --- they are not two versions
of one prediction. Its head scores the target pairs \jdg{nnRatio}$\times$ above randomly chosen
position pairs (\jdg{nnOnTarget} against \jdg{nnOffTarget}) and its output matrix is symmetric to
\jdg{nnAsym}, both checked before any contrast was computed.

\paragraph{Parameter sets are not the axis that matters.} Repeating the thermodynamic arm under
\jdg{panN} parameter sets --- \jdg{panNames}, two of them fitted to data rather than derived from
optical melting, and correlating with one another only \jdg{panCorrLo} to \jdg{panCorrHi} --- gives
skill slopes from \jdg{panLo} to \jdg{panHi}. Changing parameters within one implementation moves
the instrument's reading much less than changing model class does.

\subsection{The leakage audit, in the main text}
\label{sec:leakage}

RibonanzaNet was trained on Eterna's OpenKnot campaigns and this corpus is an OpenKnot round, so
overlap is the first thing a reader should demand and not the last. The mutants themselves cannot
be implicated --- the mutate-map-rescue library postdates the released weights --- but the wild
types can, and a memorised wild type would give the model a prior on the molecule.

We assembled \jdg{lkTrain} unique sequences from \jdg{lkDeposits} of the model's own source
libraries, deposited by the same laboratory under names that identify them, and ran four levels.
Exact whole-construct overlap with our \jdg{lkOurs} sequences is \jdg{lkExact}. At the design
window, \jdg{lkGt08} of our sequences reach $\ge0.8$ identity to some training window,
\jdg{lkGt09} reach $\ge0.9$, and \jdg{lkVerbatim} appear verbatim --- players reuse their own
designs across rounds, which is neither surprising nor benign. Excising every design cell above
the \jdg{clThr} threshold removes \jdg{clLeaked} of \jdg{clDesigns} cells and \jdg{clDirty} of
\jdg{nAdmitted} interventions. On the remaining \jdg{clClean}:

\begin{center}
\begin{tabular}{lcc}
\toprule
within-stratum skill slope & all interventions & leakage-excised \\
\midrule
RibonanzaNet & \jdg{clNeuralAll} & \jdg{clNeuralClean} \\
ViennaRNA (negative control) & \jdg{clViennaAll} & \jdg{clViennaClean} \\
\bottomrule
\end{tabular}
\end{center}

Neither skill depends on the overlap. The thermodynamic row is the control that makes this
readable: it was trained on nothing, and it behaves the same way.

\begin{scopebox}
\textbf{What the audit bounds.} These deposits are the neural model's source libraries, not a
manifest of the released checkpoint's training set. The licensed statement is \emph{no overlap
with the training sources we could obtain}, never \emph{no leakage}.
\end{scopebox}

% =====================================================================
\section{What the instrument reads out}
\label{sec:result}

Three contrasts, on the same \jdg{nAdmitted} interventions, the same stratification and the same
\jdg{resWeight} matched pair-weights. The first is the measurement. The second and third are what
each model predicts about the same pairs --- not their agreement with a reference, but the
difference their own output makes between G--C and A--U targets in matched contexts.

\begin{center}
\begin{tabular}{lccc}
\toprule
matched G--C $-$ A--U contrast in \dots & estimate & 95\% CI (puzzle-clustered) & $p$ \\
\midrule
the measurement (rescue)      & \jdg{resOut}    & \jdg{resOutCi}    & \jdg{resOutP} \\
RibonanzaNet's own prediction & \jdg{resNeural} & \jdg{resNeuralCi} & \jdg{resNeuralP} \\
ViennaRNA's own prediction    & \jdg{resVienna} & \jdg{resViennaCi} & \jdg{resViennaP} \\
\bottomrule
\end{tabular}
\end{center}

The neural model's prediction separates the two pair types by \jdg{resNeural}, in the same
direction as the measurement's \jdg{resOut} and of comparable magnitude. The thermodynamic model's
prediction shows no separation that resolves from zero. Both figures are computed twice, by
\texttt{designaudit}'s bincount estimator and by an independently written dictionary
implementation, agreeing to \jdg{resAgree}.

\begin{claimbox}
\textbf{Claim 2.} On \jdg{nAdmitted} executed two-nucleotide interventions, a neural predictor's
own output distinguishes G--C from A--U targets in matched contexts by \jdg{resNeural}
(\jdg{resNeuralCi}), tracking the measurement's \jdg{resOut} (\jdg{resOutCi}); a thermodynamic
predictor's output does not resolve such a separation (\jdg{resVienna}, \jdg{resViennaCi}). The
instrument distinguishes something one model class represents and the other does not, without a
reference structure entering the comparison.
\end{claimbox}

\paragraph{What we are not claiming.} The thermodynamic point estimate is negative, and that
negative sign holds in \jdg{resStrataNeg} of \jdg{resStrata} depth-by-crossing strata and in
\jdg{resCellsNeg} of \jdg{resCells} design cells (sign test $p=\jdg{resCellsP}$). We nevertheless
do not claim the model has the identity axis reversed. The puzzle-clustered interval covers zero,
and the design cell is not the independent unit --- cells sharing a target are not independent, so
a sign test over them overstates its own evidence. The suggestive pattern is reported here, once,
with its interval, and no conclusion rests on it.

\paragraph{An upper bound on what could run through each model.} Multiplying each predictor's own
matched contrast by its within-stratum slope on the outcome (\jdg{gmVienna} and \jdg{gmNeural})
and dividing by the measured contrast bounds the share of the identity effect that could be
carried by that predictor: at most \jdg{bdVienna} for the thermodynamic model and \jdg{bdNeural}
for the neural one. These are product-route quantities and our own linearity rule confines them to
being bounds; they are reported as such and never as point estimates.

% =====================================================================
\section{A metric we withdraw}
\label{sec:withdrawal}

\begin{withdrawbox}
\textbf{Retraction.} Earlier versions of this work reported, for each predictor, the fraction of
the identity contrast surviving after that predictor's output is held fixed by stratification, and
read values near $100\%$ as evidence that the model explains almost none of the effect.
\textbf{Those retention figures are withdrawn.} They do not measure what they were read as
measuring on this corpus, and no claim in this paper uses them.
\end{withdrawbox}

The reason is a calibration that had not been run. A retention figure is informative only if the
estimator would report a low value when mediation is genuinely present, and testing that requires
a case where the mediated fraction is known. Our first attempt at such a control was itself wrong
and is worth recording: it built a synthetic mediator depending on the treatment but not
predicting the outcome. Such a variable lies on no causal path, so conditioning on it destroys
overlap between the arms without shrinking the within-stratum contrast; the weight fell to zero
and the run reported that as detection. It had detected overlap collapse.

The second attempt changes one thing and keeps the rest real. The mediator \emph{is} the real
predictor, with its real distribution, its real correlation with design cells and its real binning
behaviour; only the outcome is synthetic,
\[
Y \;=\; \gamma\,M_{\text{real}} \;+\; \delta\,\mathbf{1}[\text{G--C}] \;+\; \varepsilon,
\]
with $\gamma$ and $\delta$ chosen so the total stratified contrast equals the corpus's own
\jdg{resOut} and a known fraction $f$ of it runs through the mediator. Sweeping $f$ says what a
measured retention means here.

\begin{center}
\begin{tabular}{lcccccc}
\toprule
true mediated fraction $f$ & $0\%$ & $10\%$ & $25\%$ & $50\%$ & $75\%$ & $100\%$ \\
\midrule
measured retention, RibonanzaNet & \jdg{mcNeuralZero} & --- & --- & --- & --- & \jdg{mcNeuralFull} \\
measured retention, ViennaRNA    & --- & --- & --- & --- & --- & \jdg{mcViennaFull} \\
\bottomrule
\end{tabular}
\end{center}

Across the sweep the neural arm's correlation with the truth is \jdg{mcNeuralCorr} --- the wrong
sign --- with mean absolute error \jdg{mcNeuralMae}; the thermodynamic arm's mean absolute error is
\jdg{mcViennaMae}, and at $f=100\%$ it returns \jdg{mcViennaFull}. At this corpus's
signal-to-noise the retention ratio is a small noisy number divided by another small noisy number,
and it carries no information about the mediated fraction.

Two things made this hard to see from the inside. The placebo mediator, which the method requires
and which we ran, controls for the mechanical shrinkage that any extra stratifying dimension
produces --- but a placebo has no relation to the outcome, so it cannot calibrate the bias that
arises precisely \emph{because} the mediator does predict the outcome. And the real predictor
separates the two arms by only \jdg{mcSep} pooled standard deviations, so there was never much for
the conditioning to remove; a value near $100\%$ was close to being forced by the design.

The quantity that replaces it is in Section~\ref{sec:result}, and it is a difference of two
matched contrasts rather than a ratio of them.

% =====================================================================
\section{Does the instrument transfer?}
\label{sec:second}

The mining criterion is stated for any probing corpus, and whether real corpora other than this
one supply interventions is a question about the world rather than about the method. We swept
\jdg{scDeposits} deposits of a public probing database under the criterion. Eleven contain
compensatory doubles, \jdg{scFlips} orientation flips in total; the rest are single-mutant
libraries, which supply none by construction.

Two of the eleven are natural RNA, and we take both forward. The first is a mutate-and-map library
of \jdg{xnFamilies} coronavirus 5$'$UTR stem-loops, measured under \jdg{xnProbes} independent
chemical probes --- dimethyl sulfate, which acylates the Watson--Crick face of adenine and
cytosine, and 2A3, which reports backbone flexibility. Neither deposit declares a target
structure, so the admission clause requiring the two positions to pair in a declared structure
cannot run and is dropped. Dropping it makes the test stricter rather than weaker: the scan is
exhaustive, so the predictor must say \emph{which} flips are rescued rather than score pairs
already known to exist. No reference structure, and no predicted structure, enters anywhere in
this section.

Under a signal-to-noise cut fixed at \jdg{xnSnCut} before the data were touched, the criterion
admits \jdg{xnAdmitDms} interventions on the first probe and \jdg{xnAdmitTwoA} on the second,
\jdg{xnShared} of them shared. Stratification here is by molecule family alone, since neither
helix depth nor crossing is defined without a structure.

\begin{claimbox}
\textbf{Both model classes retain measurable skill on natural RNA, under both probes.} On the
within-stratum rank correlation between prediction and measured rescue, the thermodynamic model
scores \jdg{xnRhoVDms} ($p=\jdg{xnRhoVDmsP}$) and \jdg{xnRhoVTwoA} ($p=\jdg{xnRhoVTwoAP}$) on the
two probes, and the neural model \jdg{xnRhoNDms} ($p=\jdg{xnRhoNDmsP}$) and \jdg{xnRhoNTwoA}
($p=\jdg{xnRhoNTwoAP}$). \textbf{The two probes are not distinguishable}: a paired cluster
bootstrap on the \emph{difference} of the two correlations returns \jdg{xnDiffV} for the
thermodynamic model and \jdg{xnDiffN} for the neural one, both covering zero. The predictors
depend only on sequence and are bit-identical across the two probes, so this is a paired
comparison in which only the measurement changes.
\end{claimbox}

\begin{table}[!htbp]
\centering\small
\caption{\textbf{The instrument on natural RNA.} Both probes measure the same interventions on the
same molecules. \emph{Rank} is the within-family rank correlation, which does not depend on the
outcome's scale; \emph{slope} is the pre-registered linear statistic, reported unchanged.}
\label{tab:transfer}
\begin{tabular}{llcccc}
\toprule
& & \multicolumn{2}{c}{rank correlation} & \multicolumn{2}{c}{slope (pre-registered)} \\
\cmidrule(lr){3-4}\cmidrule(lr){5-6}
predictor & probe & estimate & $p$ & estimate & \\
\midrule
thermodynamic & DMS & \jdg{xnRhoVDms} & \jdg{xnRhoVDmsP} & \jdg{xnSlopeVDms} & \\
              & 2A3 & \jdg{xnRhoVTwoA} & \jdg{xnRhoVTwoAP} & \jdg{xnSlopeVTwoA} & ($p=\jdg{xnSlopeVTwoAP}$) \\
neural        & DMS & \jdg{xnRhoNDms} & \jdg{xnRhoNDmsP} & \jdg{xnSlopeNDms} & \\
              & 2A3 & \jdg{xnRhoNTwoA} & \jdg{xnRhoNTwoAP} & \jdg{xnSlopeNTwoA} & \\
\bottomrule
\end{tabular}
\end{table}

\paragraph{Why the cross-probe comparison is not made on the slope.} The statistic declared in
advance was a linear within-stratum slope, and on that statistic one of the four arms --- the
thermodynamic model on 2A3 --- does not reach the threshold ($\jdg{xnSlopeVTwoA}$,
$p=\jdg{xnSlopeVTwoAP}$). That line is reported as declared. It is not, however, the statistic a
comparison \emph{between} probes may use: the two outcomes differ in dispersion by a factor of
\jdg{xnSdRatio} (\jdg{xnSdDms} against \jdg{xnSdTwoA}, with \jdg{xnOutliersDms} points below
$-1$ on the first probe and none on the second), so raw slopes are not on a common scale. The
rank statistic is, and Table~\ref{tab:transfer} reports both. That a pre-registered statistic can
be the wrong one to read is a general point about this kind of design and we take it up in a
companion paper rather than here.

\paragraph{One corpus the instrument cannot use, and why.} The second natural deposit is an
exhaustive two-dimensional double-mutant scan of three coronavirus frameshift elements ---
\jdg{xnFseFlips} orientation flips, twenty-five times what the stem-loop library supplies. It
yields \jdg{xnFseAdmit} usable interventions. The reason is not the interventions but the
references: the three wild-type profiles carry signal-to-noise between \jdg{xnFseSnLo} and
\jdg{xnFseSnHi}, against \jdg{xnSl5SnLo} and above for the stem-loop library, and the wild-type
profile is subtracted in all three difference vectors that define the outcome, so a family whose
reference fails takes every intervention in that family with it. \textbf{Whether a corpus can
carry this instrument is decided by its reference measurements, not by how many interventions it
holds} --- a constraint on corpus design that we state here and do not pursue further.

\begin{scopebox}
\textbf{Scope of this section.} It establishes that the mining step and the scoring step run on
natural molecules, under two independent probe chemistries, with no reference structure. It does
not establish an effect size transferable to OK7a: three families and one stratifying variable
give this arm far less resolution, the identity contrast here is underpowered by design, and its
interval is reported in the declaration's terms as a measurement of power rather than as evidence
about identity. Extending the transfer test to corpora with well-measured references, and to
non-viral natural RNA, is left to future work.
\end{scopebox}

\section{Discussion}
\label{sec:discussion}

\paragraph{What changes when the measured object changes.}
A structure-agreement benchmark asks whether a model's picture of a molecule is right. An
intervention benchmark asks whether the model's picture responds correctly when the molecule is
changed. The second question is harder to saturate, because it is about a difference rather than a
state, and it cannot be tuned by choosing which pairs to score, because the choice was made in the
laboratory. On this corpus the two questions give answers that differ by \jdg{cvPaired}
(\jdg{cvPairedCi}) on identical data, which is the narrow empirical version of the claim.

\paragraph{The interventions are not ours, and that is the point.}
We did not choose these mutants and we did not make them. The method's value is that it makes a
completed experiment yield an answer it was not designed to give, at no experimental cost. The
corresponding limitation is exact and stated in Section~\ref{sec:mining}: the corpus decides what
can be asked. A field that stops depositing mutant libraries stops supplying this instrument.

\paragraph{Which pairs were probed was somebody's choice.}
Mining precision is not representativeness. The experimenters chose which pairs to give a
compensatory mutant, and a criterion that admits those pairs faithfully still inherits their
selection. This is why every estimate here conditions on the design and why
Assumption~\ref{ass:placement} is stated rather than assumed silently; it is also why the
exhaustive scan noted in Section~\ref{sec:second} is scientifically more attractive than its size
suggests.

\paragraph{Two model classes are two model classes.}
We evaluated a nearest-neighbour thermodynamic implementation under \jdg{panN} parameter sets and
one transformer trained on chemical mapping. The licensed statement concerns those predictors and
those parameterisations. Nothing here supports a claim about all existing models, and the
declaration governing these runs forbade that phrase in advance.

\paragraph{Beyond RNA.}
The primitive --- recover interventions a corpus already contains, score a model on its predicted
response, condition on whoever chose the interventions --- is not specific to base pairs, and
corpora produced by design programs are becoming common in other fields. We have not tested it
elsewhere and make no claim about it; the estimator theory required for the conditioning step is
developed in a companion study, which this paper cites but does not depend on.

% =====================================================================
\section{Limitations}
\label{sec:limitations}

\begin{enumerate}[leftmargin=1.4em,itemsep=3pt]
\item \textbf{The intervention is narrow.} An orientation flip preserves the unordered pair type.
It tests whether a model's prediction of a substitution's consequences was borne out; it does not
deliver the causal effect of pair identity.
\item \textbf{Transfer is shown, its size is not.} Section~\ref{sec:second} establishes that both
steps run on natural molecules under two probe chemistries, on \jdg{xnFamilies} families and one
stratifying variable. That is enough to answer whether the instrument reads, and not enough to
carry an effect size across corpora; the second corpus is a demonstration of scope, not a
replication. A corpus whose reference profiles are poorly measured yields nothing at all, as
\jdg{xnFseFlips} flips reduced to \jdg{xnFseAdmit} usable ones show.
\item \textbf{The leakage audit bounds what it can reach.} \jdg{lkTrain} sequences from the neural
model's source libraries are not the released checkpoint's training manifest.
\item \textbf{A negative result we cannot resolve.} The thermodynamic model's own matched contrast
is negative in point estimate and in most strata, but its clustered interval covers zero. Whether
that model is merely insensitive on this axis or actively inverted is not settled here.
\item \textbf{A retracted metric.} Section~\ref{sec:withdrawal} withdraws the retention figures
this work previously reported. The design that produced them is not repaired by this paper; it is
replaced by a quantity that does not require it.
\end{enumerate}

% =====================================================================
\section{Conclusion}

Compensatory mutagenesis in a chemical-probing corpus is a controlled two-nucleotide intervention
that has already been performed \jdg{nAdmitted} times. Recovering those interventions by a stated
criterion, and scoring predictors on whether their output responds as the measurement did, gives a
readout that is not saturated where structure agreement is, that admits no choice of which pairs
to score, and that separates a thermodynamic from a neural predictor on an axis structure
agreement cannot express. The instrument was calibrated before it was used, audited for leakage in
the main text, and used to retract one of this work's own earlier numbers. It costs no new
experiments, and it runs on data that is already public.

% =====================================================================
\bibliographystyle{unsrtnat}
\bibliography{refs}

\appendix
% =====================================================================
\section{Preregistration record}
\label{app:prereg}

Every experiment in this paper was declared in a frozen JSON file committed and pushed before the
code that produced it existed, naming the arms, the kill conditions, the probes that must fire,
and a list of staked predictions written before the run. The table below records the outcome of
each stake. Predictions that failed are marked as failed; the accompanying text does not
re-describe them as successes.

\begin{center}
\begin{tabular}{p{0.62\textwidth}c}
\toprule
staked before the run & outcome \\
\midrule
Neural model has positive skill on this instrument, $p\le0.005$ & correct \\
Neural skill exceeds the thermodynamic model's & correct \\
Neural calibration is monotone & \textbf{failed} \\
Two predictors correlate below the same-class threshold & correct \\
Exact whole-sequence leakage is zero & correct \\
Leaked design cells are at most $10\%$ of the corpus & correct \\
Clean-subset skill differs from leaked-subset skill by less than $0.05$ & \textbf{failed} \\
Marginal estimator bias equals $\beta g$ & correct \\
Matching is necessary to remove placement bias & \textbf{failed} \\
\bottomrule
\end{tabular}
\end{center}

\section{The criterion as code}
\label{app:code}

\texttt{designaudit mine} implements Definition~\ref{def:flip} exactly, with each exclusion
reason counted and returned. \texttt{designaudit score} returns the three readouts of
Section~\ref{sec:scoring} and refuses to report a residual when sensitivity is null.
\texttt{designaudit selftest} reproduces \jdg{nAdmitted}, \jdg{resOut}, \jdg{resWeight} and
\jdg{skVienna} from the public corpus file. Predictors are supplied as
\texttt{module:callable} taking a sequence and two positions.

\end{document}
